\section{環境法・循環・持続可能性}

\subsection{平成29年度 問5-1}
\label{q:h29-reg-basic-5-1}

\noindent\textbf{原典：} sources/H29/basic.pdf、PDF物理18頁。

\textbf{公式正答：} 1（single）。

\par\textit{原典本文の校正と独自解説は未作成。分類のみの登録項目。}

\subsection{平成30年度 問5-1}
\label{q:h30-reg-basic-5-1}

\noindent\textbf{原典：} sources/H30/basic.pdf、PDF物理18頁。

\textbf{公式正答：} 2（single）。

\par\textit{原典本文の校正と独自解説は未作成。分類のみの登録項目。}

\subsection{平成30年度 問5-2}
\label{q:h30-reg-basic-5-2}

\noindent\textbf{原典：} sources/H30/basic.pdf、PDF物理18頁。

\textbf{公式正答：} 3（single）。

\par\textit{原典本文の校正と独自解説は未作成。分類のみの登録項目。}

\subsection{令和元年度 問5-1}
\label{q:r01-reg-basic-5-1}

\noindent\textbf{原典：} sources/R01/basic.pdf、PDF物理19頁。

\textbf{公式正答：} 3（single）。

\par\textit{原典本文の校正と独自解説は未作成。分類のみの登録項目。}

\subsection{令和元年度 問5-2}
\label{q:r01-reg-basic-5-2}

\noindent\textbf{原典：} sources/R01/basic.pdf、PDF物理20頁。

\textbf{公式正答：} 2（single）。

\par\textit{原典本文の校正と独自解説は未作成。分類のみの登録項目。}

\subsection{令和元年度再試験 問5-1}
\label{q:r01-ret-basic-5-1}

\noindent\textbf{原典：} sources/R01\_retest/basic.pdf、PDF物理17頁。

\textbf{公式正答：} 1（single）。

\par\textit{原典本文の校正と独自解説は未作成。分類のみの登録項目。}

\subsection{令和元年度再試験 問5-2}
\label{q:r01-ret-basic-5-2}

\noindent\textbf{原典：} sources/R01\_retest/basic.pdf、PDF物理18頁。

\textbf{公式正答：} 1（single）。

\par\textit{原典本文の校正と独自解説は未作成。分類のみの登録項目。}

\subsection{令和元年度再試験 問5-4}
\label{q:r01-ret-basic-5-4}

\noindent\textbf{原典：} sources/R01\_retest/basic.pdf、PDF物理19頁。

\textbf{公式正答：} 3（single）。

\par\textit{原典本文の校正と独自解説は未作成。分類のみの登録項目。}

\subsection{令和2年度 問5-1}
\label{q:r02-reg-basic-5-1}

\noindent\textbf{原典：} sources/R02/basic.pdf、PDF物理19頁。

\textbf{公式正答：} 2（single）。

\par\textit{原典本文の校正と独自解説は未作成。分類のみの登録項目。}

\subsection{令和3年度 問5-2}
\label{q:r03-reg-basic-5-2}

\noindent\textbf{原典：} sources/R03/basic.pdf、PDF物理22頁。

\textbf{公式正答：} 5（single）。

\par\textit{原典本文の校正と独自解説は未作成。分類のみの登録項目。}

\subsection{令和3年度 問5-3}
\label{q:r03-reg-basic-5-3}

\noindent\textbf{原典：} sources/R03/basic.pdf、PDF物理23頁。

\textbf{公式正答：} 1（single）。

\par\textit{原典本文の校正と独自解説は未作成。分類のみの登録項目。}

\subsection{令和4年度 問5-2}
\label{q:r04-reg-basic-5-2}

\noindent\textbf{原典：} sources/R04/basic.pdf、PDF物理17頁。

\textbf{公式正答：} 1（single）。

\par\textit{原典本文の校正と独自解説は未作成。分類のみの登録項目。}

\subsection{令和6年度 問5-1}
\label{q:r06-reg-basic-5-1}

\noindent\textbf{原典：} sources/R06/basic.pdf、PDF物理17頁。

\textbf{公式正答：} 1（single）。

\par\textit{原典本文の校正と独自解説は未作成。分類のみの登録項目。}

\subsection{令和6年度 問5-2}
\label{q:r06-reg-basic-5-2}

\noindent\textbf{原典：} sources/R06/basic.pdf、PDF物理18頁。

\textbf{公式正答：} 5（single）。

\par\textit{原典本文の校正と独自解説は未作成。分類のみの登録項目。}

\subsection{令和6年度 問5-3}
\label{q:r06-reg-basic-5-3}

\noindent\textbf{原典：} sources/R06/basic.pdf、PDF物理19頁。

\textbf{公式正答：} 2（single）。

\par\textit{原典本文の校正と独自解説は未作成。分類のみの登録項目。}

\subsection{令和6年度 問5-4}
\label{q:r06-reg-basic-5-4}

\noindent\textbf{原典：} sources/R06/basic.pdf、PDF物理19頁。

\textbf{公式正答：} 4（single）。

\par\textit{原典本文の校正と独自解説は未作成。分類のみの登録項目。}

\subsection{令和7年度 問5-1}
\label{q:r07-reg-basic-5-1}

\noindent\textbf{原典：} sources/R07/basic.pdf、PDF物理15頁。

\textbf{公式正答：} 5（single）。

\par\textit{原典本文の校正と独自解説は未作成。分類のみの登録項目。}

\subsection{令和7年度 問5-2}
\label{q:r07-reg-basic-5-2}

\noindent\textbf{原典：} sources/R07/basic.pdf、PDF物理16頁。

\textbf{公式正答：} 3（single）。

\par\textit{原典本文の校正と独自解説は未作成。分類のみの登録項目。}

\subsection{令和7年度 問5-4}
\label{q:r07-reg-basic-5-4}

\noindent\textbf{原典：} sources/R07/basic.pdf、PDF物理18頁。

\textbf{公式正答：} 2（single）。

\par\textit{原典本文の校正と独自解説は未作成。分類のみの登録項目。}

\section{科学技術史}

\subsection{平成29年度 問5-5}
\label{q:h29-reg-basic-5-5}

\noindent\textbf{原典：} sources/H29/basic.pdf、PDF物理19頁。

\textbf{公式正答：} 2（single）。

\par\textit{原典本文の校正と独自解説は未作成。分類のみの登録項目。}

\subsection{令和6年度 問5-5}
\label{q:r06-reg-basic-5-5}

\noindent\textbf{原典：} sources/R06/basic.pdf、PDF物理20頁。

\textbf{公式正答：} 3（single）。

\par\textit{原典本文の校正と独自解説は未作成。分類のみの登録項目。}

\section{知的財産}

\subsection{令和元年度 問5-6}
\label{q:r01-reg-basic-5-6}

\noindent\textbf{原典：} sources/R01/basic.pdf、PDF物理22頁。

\textbf{公式正答：} 5（single）。

\par\textit{原典本文の校正と独自解説は未作成。分類のみの登録項目。}

\section{安全・リスク評価}

\subsection{平成30年度 問5-3}
\label{q:h30-reg-basic-5-3}

\noindent\textbf{原典：} sources/H30/basic.pdf、PDF物理19頁。

\textbf{公式正答：} 3（single）。

\par\textit{原典本文の校正と独自解説は未作成。分類のみの登録項目。}

\subsection{令和元年度再試験 問5-6}
\label{q:r01-ret-basic-5-6}

\noindent\textbf{原典：} sources/R01\_retest/basic.pdf、PDF物理20頁。

\textbf{公式正答：} 3（single）。

\par\textit{原典本文の校正と独自解説は未作成。分類のみの登録項目。}

\subsection{令和4年度 問5-3}
\label{q:r04-reg-basic-5-3}

\noindent\textbf{原典：} sources/R04/basic.pdf、PDF物理18頁。

\textbf{公式正答：} 1（single）。

\par\textit{原典本文の校正と独自解説は未作成。分類のみの登録項目。}

\subsection{令和4年度 問5-5}
\label{q:r04-reg-basic-5-5}

\noindent\textbf{原典：} sources/R04/basic.pdf、PDF物理20頁。

\textbf{公式正答：} 5（single）。

\par\textit{原典本文の校正と独自解説は未作成。分類のみの登録項目。}

\section{環境会計と評価・報告・監査を区別}

\subsection{平成26年度 問5-2}
\label{q:h26-reg-basic-5-2}

\noindent\textbf{原典：} sources/H26/basic.pdf、PDF物理16頁。

\textbf{公式正答：} 5（single）。

\paragraph{問題の要旨} 環境会計と評価・報告・監査を区別。

\Needspace{4\baselineskip}\paragraph{確認した条件・記述（独自要約）}\mbox{}\par

\noindent\textbf{時点：} 平成26年度に指定された資料・制度。\par

\noindent\textbf{照合：} 原典21画像で否定・五肢・組合せ表を確認。後年値で置換しない。\par

\paragraph{問われる核心} 環境会計と評価・報告・監査を区別。

\paragraph{必要知識} 対象、評価目的、会計単位、環境管理手段。

\paragraph{考え方・正答を選ぶ理由（査読待ち初稿）} 環境会計はコストと効果を認識して、貨幣単位又は物量単位で定量化・伝達する仕組みで5。LCAは設備の更新周期だけでなく製品等の全過程。環境報告は社会への情報開示、グリーン購入は必要性と環境負荷、環境監査は環境管理の点検であり、税務監査だけではない。

\Needspace{4\baselineskip}\paragraph{選択肢の検討（査読待ち初稿）}\mbox{}\par

\noindent\textbf{1：不適切} LCAは設備更新時のみ。式・定義・対応の独立判定と一致しない。\par

\noindent\textbf{2：不適切} 環境報告書は設備の自治体届出。式・定義・対応の独立判定と一致しない。\par

\noindent\textbf{3：不適切} グリーン購入は生物資源優先のみ。式・定義・対応の独立判定と一致しない。\par

\noindent\textbf{4：不適切} 環境監査は税務処理のみ。式・定義・対応の独立判定と一致しない。\par

\noindent\textbf{5：適切・選択対象} 環境会計は費用効果を金額物量で。判定と一致する。\par

\paragraph{再利用する考え方} 対象、評価目的、会計単位、環境管理手段。 指定年度・対象範囲・単位・否定を先に固定し、他年度の番号や現行値を転用しない。

\paragraph{共通理論} \hyperref[theory:env:accounting]{第\ref*{theory:env:accounting}節}

\par 参考：\cite{env58_accounting}。

\par 難易度 2（校準初稿）、優先度未評価。

\section{粒径範囲と大気汚染の生成経路}

\subsection{平成26年度 問5-1}
\label{q:h26-reg-basic-5-1}

\noindent\textbf{原典：} sources/H26/basic.pdf、PDF物理16頁。

\textbf{公式正答：} 1（single）。

\par\textbf{検討未確定：} COによるメタン寿命の化学過程と各健康影響は追加一次資料照合。肺沈着と粒径の関係を2.5μm一値へ固定しない。

\paragraph{問題の要旨} 粒径範囲と大気汚染の生成経路。

\Needspace{4\baselineskip}\paragraph{確認した条件・記述（独自要約）}\mbox{}\par

\noindent\textbf{時点：} 平成26年度に指定された資料・制度。\par

\noindent\textbf{照合：} 原典21画像で否定・五肢・組合せ表を確認。後年値で置換しない。\par

\paragraph{問われる核心} 粒径範囲と大気汚染の生成経路。

\paragraph{必要知識} 大きさの範囲と測定、一次・二次汚染物質、生成と除去。

\paragraph{考え方・正答を選ぶ理由（査読待ち初稿）} PM2.5は概ね粒径2.5μm以下の微小粒子で、2.5付近だけに限定する1が誤り。規格上は2.5μmを50\%分離する装置で採取するという測定定義もある。SO2は硫黄、NOxとVOCは二次生成、COは不完全燃焼とヘモグロビンの結合。COによるOH消費はメタン除去に影響する。

\Needspace{4\baselineskip}\paragraph{選択肢の検討（査読待ち初稿）}\mbox{}\par

\noindent\textbf{1：不適切・選択対象} PM2.5は2.5μm付近だけ。判定と一致する。\par

\noindent\textbf{2：条件内で適切} SO2は燃料中S由来。式・定義・対応の独立判定と一致しない。\par

\noindent\textbf{3：条件内で適切} NO2は燃焼と光化学前駆物質。式・定義・対応の独立判定と一致しない。\par

\noindent\textbf{4：条件内で適切} 光化学オキシダントはNOxとVOCの光反応。式・定義・対応の独立判定と一致しない。\par

\noindent\textbf{5：条件内で適切} COは不完全燃焼でメタン寿命にも影響。式・定義・対応の独立判定と一致しない。\par

\paragraph{再利用する考え方} 大きさの範囲と測定、一次・二次汚染物質、生成と除去。 指定年度・対象範囲・単位・否定を先に固定し、他年度の番号や現行値を転用しない。

\paragraph{共通理論} \hyperref[theory:env:air]{第\ref*{theory:env:air}節}

\par 参考：\cite{env58_pmstandard,env58_pm25}。

\par 難易度 3（校準初稿）、優先度未評価。

\section{産業廃棄物フローを物質収支で検算}

\subsection{平成28年度 問5-1}
\label{q:h28-reg-basic-5-1}

\noindent\textbf{原典：} sources/H28/basic.pdf、PDF物理21頁。

\textbf{公式正答：} 5（single）。

\paragraph{問題の要旨} 産業廃棄物フローを物質収支で検算。

\Needspace{4\baselineskip}\paragraph{確認した条件・記述（独自要約）}\mbox{}\par

\noindent\textbf{排出：} 379百万t。\par

\noindent\textbf{中間処理：} 290。\par

\noindent\textbf{直接再生：} 83。\par

\noindent\textbf{処理後再生：} 124。\par

\noindent\textbf{処理後最終処分：} 7。\par

\noindent\textbf{全再生：} 207。\par

\noindent\textbf{全最終処分：} 13。\par

\paragraph{問われる核心} 産業廃棄物フローを物質収支で検算。

\paragraph{必要知識} 分岐の収支、直接と処理後、百分率と比率。

\paragraph{考え方・正答を選ぶ理由（査読待ち初稿）} 単位百万t。中間処理290は再利用124、処理後処分7、減量159へ分かれる。総処分13から直接処分は13−7=\allowbreak{}6で、60ではないので5。総排出も83+290+6=\allowbreak{}379、最終の207+159+13=\allowbreak{}379で一致。83/207≈.40、207/379≈.55、207/13≈16、124/83≈1.5、159/379≈.42。

\Needspace{4\baselineskip}\paragraph{選択肢の検討（査読待ち初稿）}\mbox{}\par

\noindent\textbf{1：条件内で適切} 直接再利用83で約40\%。式・定義・対応の独立判定と一致しない。\par

\noindent\textbf{2：条件内で適切} 再利用207で55\%・処分の16倍。式・定義・対応の独立判定と一致しない。\par

\noindent\textbf{3：条件内で適切} 処理後124で直接の1.5倍。式・定義・対応の独立判定と一致しない。\par

\noindent\textbf{4：条件内で適切} 減量159で42\%。式・定義・対応の独立判定と一致しない。\par

\noindent\textbf{5：不適切・選択対象} 直接処分60で16\%。判定と一致する。\par

\paragraph{再利用する考え方} 分岐の収支、直接と処理後、百分率と比率。 指定年度・対象範囲・単位・否定を先に固定し、他年度の番号や現行値を転用しない。

\paragraph{共通理論} \hyperref[theory:env:balance]{第\ref*{theory:env:balance}節}

\par 難易度 2（校準初稿）、優先度未評価。

\section{生物多様性条約と外来種防除}

\subsection{平成28年度 問5-2}
\label{q:h28-reg-basic-5-2}

\noindent\textbf{原典：} sources/H28/basic.pdf、PDF物理22頁。

\textbf{公式正答：} 5（single）。

\par\textbf{検討未確定：} 1992・1993と議定書、当時の外来生物法の防除条項の一次本文を追加照合。動物愛護上の配慮が不要という意味にはしない。

\paragraph{問題の要旨} 生物多様性条約と外来種防除。

\Needspace{4\baselineskip}\paragraph{確認した条件・記述（独自要約）}\mbox{}\par

\noindent\textbf{時点：} 平成28年度に指定された資料・制度。\par

\noindent\textbf{照合：} 原典21画像で否定・五肢・組合せ表を確認。後年値で置換しない。\par

\paragraph{問われる核心} 生物多様性条約と外来種防除。

\paragraph{必要知識} 条約目的、署名と発効、外来種と防除。

\paragraph{考え方・正答を選ぶ理由（査読待ち初稿）} 外来種による生態系被害の防止には防除制度があり、一律に駆除を禁止する5が誤り。CBDの三目的は保全、持続可能な利用、遺伝資源の利益の公正衡平配分。組換え生物の扱いはCartagena議定書であり、遺伝資源の利益配分のNagoya議定書と混同しない。

\Needspace{4\baselineskip}\paragraph{選択肢の検討（査読待ち初稿）}\mbox{}\par

\noindent\textbf{1：条件内で適切} 組換え生物の議定書。式・定義・対応の独立判定と一致しない。\par

\noindent\textbf{2：条件内で適切} 1992署名開放・1993発効。式・定義・対応の独立判定と一致しない。\par

\noindent\textbf{3：条件内で適切} 保全・持続利用・利益配分。式・定義・対応の独立判定と一致しない。\par

\noindent\textbf{4：条件内で適切} 外来種に重大影響の例。式・定義・対応の独立判定と一致しない。\par

\noindent\textbf{5：不適切・選択対象} 日本は動物愛護で駆除禁止。判定と一致する。\par

\paragraph{再利用する考え方} 条約目的、署名と発効、外来種と防除。 指定年度・対象範囲・単位・否定を先に固定し、他年度の番号や現行値を転用しない。

\paragraph{共通理論} \hyperref[theory:env:biodiversity]{第\ref*{theory:env:biodiversity}節}

\par 難易度 3（校準初稿）、優先度未評価。

\section{2008年度の温室効果ガスと国際枠組}

\subsection{平成23年度 問5-1}
\label{q:h23-reg-basic-5-1}

\noindent\textbf{原典：} sources/H23/basic.pdf、PDF物理15頁。

\textbf{公式正答：} 1（single）。

\par\textbf{検討未確定：} 当時白書の全ガス内訳・部門図と2009合意の条文の追加照合が必要。CO2九割以上は当時確定値公表表の検索本文でも確認。

\paragraph{問題の要旨} 2008年度の温室効果ガスと国際枠組。

\Needspace{4\baselineskip}\paragraph{確認した条件・記述（独自要約）}\mbox{}\par

\noindent\textbf{時点：} 平成23年度に指定された資料・制度。\par

\noindent\textbf{照合：} 原典21画像で否定・五肢・組合せ表を確認。後年値で置換しない。\par

\paragraph{問われる核心} 2008年度の温室効果ガスと国際枠組。

\paragraph{必要知識} CO2換算、部門の定義、条約と議定書、採択年。

\paragraph{考え方・正答を選ぶ理由（査読待ち初稿）} CO2換算の2008年度総量の九割以上がCO2で、1。部門別では家庭が第二位という組合せが違う。UNFCCCの究極目的に2050という一律期限はなく、1997採択は京都議定書。2009コペンハーゲンの合意を全途上国に法的数量削減義務を課す議定書と扱わない。

\Needspace{4\baselineskip}\paragraph{選択肢の検討（査読待ち初稿）}\mbox{}\par

\noindent\textbf{1：適切・選択対象} CO2が九割以上。判定と一致する。\par

\noindent\textbf{2：不適切} 産業の次は家庭。式・定義・対応の独立判定と一致しない。\par

\noindent\textbf{3：不適切} 条約の濃度安定は2050期限。式・定義・対応の独立判定と一致しない。\par

\noindent\textbf{4：不適切} 1997パリ議定書。式・定義・対応の独立判定と一致しない。\par

\noindent\textbf{5：不適切} 2009途上国に基準年削減数値義務。式・定義・対応の独立判定と一致しない。\par

\paragraph{再利用する考え方} CO2換算、部門の定義、条約と議定書、採択年。 指定年度・対象範囲・単位・否定を先に固定し、他年度の番号や現行値を転用しない。

\paragraph{共通理論} \hyperref[theory:env:climate]{第\ref*{theory:env:climate}節}

\par 参考：\cite{env58_ghg2008}。

\par 難易度 3（校準初稿）、優先度未評価。

\section{暖房と自動車のCO2排出量を比較}

\subsection{平成25年度 問5-2}
\label{q:h25-reg-basic-5-2}

\noindent\textbf{原典：} sources/H25/basic.pdf、PDF物理19頁。

\textbf{公式正答：} 2（single）。

\paragraph{問題の要旨} 暖房と自動車のCO2排出量を比較。

\Needspace{4\baselineskip}\paragraph{確認した条件・記述（独自要約）}\mbox{}\par

\noindent\textbf{電力：} 500W、3h/日×50日、.42kg-CO2/kWh。\par

\noindent\textbf{車：} 10km/L、200km、2.32kg-CO2/L。\par

\paragraph{問われる核心} 暖房と自動車のCO2排出量を比較。

\paragraph{必要知識} エネルギー=電力×時間、燃料使用量、排出係数。

\paragraph{考え方・正答を選ぶ理由（査読待ち初稿）} 暖房は.5kW×3h×50=\allowbreak{}75kWh、排出31.5kg。自動車は200/10=\allowbreak{}20L、排出46.4kg。比31.5/46.4≈.679で2:3に最も近い。WをkWへ直してから時間を掛ける。

\Needspace{4\baselineskip}\paragraph{選択肢の検討（査読待ち初稿）}\mbox{}\par

\noindent\textbf{1：不適切} 1:3。式・定義・対応の独立判定と一致しない。\par

\noindent\textbf{2：適切・選択対象} 2:3。判定と一致する。\par

\noindent\textbf{3：不適切} 1:1。式・定義・対応の独立判定と一致しない。\par

\noindent\textbf{4：不適切} 3:2。式・定義・対応の独立判定と一致しない。\par

\noindent\textbf{5：不適切} 3:1。式・定義・対応の独立判定と一致しない。\par

\paragraph{再利用する考え方} エネルギー=電力×時間、燃料使用量、排出係数。 指定年度・対象範囲・単位・否定を先に固定し、他年度の番号や現行値を転用しない。

\paragraph{共通理論} \hyperref[theory:env:co2]{第\ref*{theory:env:co2}節}

\par 難易度 2（校準初稿）、優先度未評価。

\subsection{平成27年度 問5-3}
\label{q:h27-reg-basic-5-3}

\noindent\textbf{原典：} sources/H27/basic.pdf、PDF物理16頁。

\textbf{公式正答：} 5（single）。

\paragraph{問題の要旨} 利用人数と時間を含めた一人当たり排出。

\Needspace{4\baselineskip}\paragraph{確認した条件・記述（独自要約）}\mbox{}\par

\noindent\textbf{a：} 10km/L、100km、3人、2.32kg/L。\par

\noindent\textbf{b：} 8000kW、200km/h、500km、1000人、.47kg/kWh。\par

\noindent\textbf{c：} 500W、4h/日、50日、4人、.47kg/kWh。\par

\paragraph{問われる核心} 利用人数と時間を含めた一人当たり排出。

\paragraph{必要知識} 排出係数、乗車人数、走行時間、総量と原単位。

\paragraph{考え方・正答を選ぶ理由（査読待ち初稿）} aは100/10×2.32/3=\allowbreak{}7.73kg/人。bは500/200=\allowbreak{}2.5h、8000×2.5×.47/1000=\allowbreak{}9.4kg/人。cは.5×4×50×.47/4=\allowbreak{}11.75kg/人。c>b>aの5。総量の大小ではなく人数で割った量を比べる。

\Needspace{4\baselineskip}\paragraph{選択肢の検討（査読待ち初稿）}\mbox{}\par

\noindent\textbf{1：不適切} a>b>c。式・定義・対応の独立判定と一致しない。\par

\noindent\textbf{2：不適切} a>c>b。式・定義・対応の独立判定と一致しない。\par

\noindent\textbf{3：不適切} b>a>c。式・定義・対応の独立判定と一致しない。\par

\noindent\textbf{4：不適切} c>a>b。式・定義・対応の独立判定と一致しない。\par

\noindent\textbf{5：適切・選択対象} c>b>a。判定と一致する。\par

\paragraph{再利用する考え方} 排出係数、乗車人数、走行時間、総量と原単位。 指定年度・対象範囲・単位・否定を先に固定し、他年度の番号や現行値を転用しない。

\paragraph{共通理論} \hyperref[theory:env:co2]{第\ref*{theory:env:co2}節}

\par 難易度 2（校準初稿）、優先度未評価。

\section{科学と社会をつなぐ双方向の対話}

\subsection{平成25年度 問5-5}
\label{q:h25-reg-basic-5-5}

\noindent\textbf{原典：} sources/H25/basic.pdf、PDF物理22頁。

\textbf{公式正答：} 2（single）。

\paragraph{問題の要旨} 科学と社会をつなぐ双方向の対話。

\Needspace{4\baselineskip}\paragraph{確認した条件・記述（独自要約）}\mbox{}\par

\noindent\textbf{四文の順序：} ア=社会理解、イ=科学技術間だけ、ウ=双方向、エ=媒体へ技術使用\par

\noindent\textbf{問：} 適切な二文。\par

\paragraph{問われる核心} 科学と社会をつなぐ双方向の対話。

\paragraph{必要知識} 専門家と社会、理解と対話、情報発信と相互学習。

\paragraph{考え方・正答を選ぶ理由（査読待ち初稿）} 社会・国民の期待や考えを専門家自身が理解することと、専門家・非専門家間の双方向性が適切。基礎科学と応用技術の内部連携だけをこの領域全体の定義にはできず、媒体制作へ技術を使うことや科学者へ任せることも、社会と対話するという目的の代替ではない。原典の列順で組合せを選ぶ。

\Needspace{4\baselineskip}\paragraph{選択肢の検討（査読待ち初稿）}\mbox{}\par

\noindent\textbf{1：不適切} ア イ。式・定義・対応の独立判定と一致しない。\par

\noindent\textbf{2：適切・選択対象} ア ウ。判定と一致する。\par

\noindent\textbf{3：不適切} イ ウ。式・定義・対応の独立判定と一致しない。\par

\noindent\textbf{4：不適切} イ エ。式・定義・対応の独立判定と一致しない。\par

\noindent\textbf{5：不適切} ウ エ。式・定義・対応の独立判定と一致しない。\par

\paragraph{再利用する考え方} 専門家と社会、理解と対話、情報発信と相互学習。 指定年度・対象範囲・単位・否定を先に固定し、他年度の番号や現行値を転用しない。

\paragraph{共通理論} \hyperref[theory:env:communication]{第\ref*{theory:env:communication}節}

\par 参考：\cite{env58_communication}。

\par 難易度 2（校準初稿）、優先度未評価。

\subsection{平成28年度 問5-5}
\label{q:h28-reg-basic-5-5}

\noindent\textbf{原典：} sources/H28/basic.pdf、PDF物理24頁。

\textbf{公式正答：} 5（single）。

\paragraph{問題の要旨} 科学と社会をつなぐ双方向の対話。

\Needspace{4\baselineskip}\paragraph{確認した条件・記述（独自要約）}\mbox{}\par

\noindent\textbf{四文の順序：} ア=科学技術間だけ、イ=科学者に制作委任、ウ=社会理解、エ=双方向\par

\noindent\textbf{問：} 適切な二文。\par

\paragraph{問われる核心} 科学と社会をつなぐ双方向の対話。

\paragraph{必要知識} 専門家と社会、理解と対話、情報発信と相互学習。

\paragraph{考え方・正答を選ぶ理由（査読待ち初稿）} 社会・国民の期待や考えを専門家自身が理解することと、専門家・非専門家間の双方向性が適切。基礎科学と応用技術の内部連携だけをこの領域全体の定義にはできず、媒体制作へ技術を使うことや科学者へ任せることも、社会と対話するという目的の代替ではない。原典の列順で組合せを選ぶ。

\Needspace{4\baselineskip}\paragraph{選択肢の検討（査読待ち初稿）}\mbox{}\par

\noindent\textbf{1：不適切} ア イ。式・定義・対応の独立判定と一致しない。\par

\noindent\textbf{2：不適切} ア ウ。式・定義・対応の独立判定と一致しない。\par

\noindent\textbf{3：不適切} イ ウ。式・定義・対応の独立判定と一致しない。\par

\noindent\textbf{4：不適切} イ エ。式・定義・対応の独立判定と一致しない。\par

\noindent\textbf{5：適切・選択対象} ウ エ。判定と一致する。\par

\paragraph{再利用する考え方} 専門家と社会、理解と対話、情報発信と相互学習。 指定年度・対象範囲・単位・否定を先に固定し、他年度の番号や現行値を転用しない。

\paragraph{共通理論} \hyperref[theory:env:communication]{第\ref*{theory:env:communication}節}

\par 参考：\cite{env58_communication}。

\par 難易度 2（校準初稿）、優先度未評価。

\section{2010年の電源比とガス技術を分ける}

\subsection{平成24年度 問5-3}
\label{q:h24-reg-basic-5-3}

\noindent\textbf{原典：} sources/H24/basic.pdf、PDF物理16頁。

\textbf{公式正答：} 2（single）。

\par\textbf{検討未確定：} 一般電気事業用という分母の2010年度原表と2005以降の米国生産統計は追加取得。現行の全発電量比率と混同しない。

\paragraph{問題の要旨} 2010年の電源比とガス技術を分ける。

\Needspace{4\baselineskip}\paragraph{確認した条件・記述（独自要約）}\mbox{}\par

\noindent\textbf{時点：} 平成24年度に指定された資料・制度。\par

\noindent\textbf{照合：} 原典21画像で否定・五肢・組合せ表を確認。後年値で置換しない。\par

\paragraph{問われる核心} 2010年の電源比とガス技術を分ける。

\paragraph{必要知識} 歴史的電源構成、複合サイクル、非在来型天然ガス。

\paragraph{考え方・正答を選ぶ理由（査読待ち初稿）} 2010年度の原子力比率は約30\%。天然ガスのコンバインドサイクルはガスタービンと排熱による蒸気タービンを組み合わせる。米国で増産が注目された非在来ガスはシェールガス。三語の組合せ2。

\Needspace{4\baselineskip}\paragraph{選択肢の検討（査読待ち初稿）}\mbox{}\par

\noindent\textbf{1：不適切} 30/再熱再生/タイトサンド。式・定義・対応の独立判定と一致しない。\par

\noindent\textbf{2：適切・選択対象} 30/コンバインド/シェール。判定と一致する。\par

\noindent\textbf{3：不適切} 30/コンバインド/タイトサンド。式・定義・対応の独立判定と一致しない。\par

\noindent\textbf{4：不適切} 20/再熱再生/シェール。式・定義・対応の独立判定と一致しない。\par

\noindent\textbf{5：不適切} 20/コンバインド/シェール。式・定義・対応の独立判定と一致しない。\par

\paragraph{再利用する考え方} 歴史的電源構成、複合サイクル、非在来型天然ガス。 指定年度・対象範囲・単位・否定を先に固定し、他年度の番号や現行値を転用しない。

\paragraph{共通理論} \hyperref[theory:env:energy2010]{第\ref*{theory:env:energy2010}節}

\par 難易度 3（校準初稿）、優先度未評価。

\section{2011年一人当たり一次エネルギー}

\subsection{平成27年度 問5-4}
\label{q:h27-reg-basic-5-4}

\noindent\textbf{原典：} sources/H27/basic.pdf、PDF物理17頁。

\textbf{公式正答：} 5（single）。

\par\textbf{検討未確定：} IEAの指定2011年国別原表は未取得。現行データページを当時値の照合証跡にはせず、Astraの追加確認対象に残す。

\paragraph{問題の要旨} 2011年一人当たり一次エネルギー。

\Needspace{4\baselineskip}\paragraph{確認した条件・記述（独自要約）}\mbox{}\par

\noindent\textbf{時点：} 平成27年度に指定された資料・制度。\par

\noindent\textbf{照合：} 原典21画像で否定・五肢・組合せ表を確認。後年値で置換しない。\par

\paragraph{問われる核心} 2011年一人当たり一次エネルギー。

\paragraph{必要知識} 一次エネルギー、石油換算、人口当たり、年次比較。

\paragraph{考え方・正答を選ぶ理由（査読待ち初稿）} 原典の2011年指標では米加が7toe以上、韓露約5toe、独日約4toeで5。1toe≈42GJを使い、国全体の総量と人口当たりを混同しない。世界1.9、中国2.0という同じ年・同じ指標で比較する。

\Needspace{4\baselineskip}\paragraph{選択肢の検討（査読待ち初稿）}\mbox{}\par

\noindent\textbf{1：不適切} 独日/米加/韓露。式・定義・対応の独立判定と一致しない。\par

\noindent\textbf{2：不適切} 韓露/独日/米加。式・定義・対応の独立判定と一致しない。\par

\noindent\textbf{3：不適切} 韓露/米加/独日。式・定義・対応の独立判定と一致しない。\par

\noindent\textbf{4：不適切} 米加/独日/韓露。式・定義・対応の独立判定と一致しない。\par

\noindent\textbf{5：適切・選択対象} 米加/韓露/独日。判定と一致する。\par

\paragraph{再利用する考え方} 一次エネルギー、石油換算、人口当たり、年次比較。 指定年度・対象範囲・単位・否定を先に固定し、他年度の番号や現行値を転用しない。

\paragraph{共通理論} \hyperref[theory:env:energycapita]{第\ref*{theory:env:energycapita}節}

\par 難易度 3（校準初稿）、優先度未評価。

\section{2015年策定の2030電源見通しを読む}

\subsection{平成28年度 問5-4}
\label{q:h28-reg-basic-5-4}

\noindent\textbf{原典：} sources/H28/basic.pdf、PDF物理23頁。

\textbf{公式正答：} 1（single）。

\par\textbf{検討未確定：} 2015公開関連資料の検索と白書骨子を参照。策定7月の最終PDFの全図表照合は未了で後工程に残す。

\paragraph{問題の要旨} 2015年策定の2030電源見通しを読む。

\Needspace{4\baselineskip}\paragraph{確認した条件・記述（独自要約）}\mbox{}\par

\noindent\textbf{時点：} 平成28年度に指定された資料・制度。\par

\noindent\textbf{照合：} 原典21画像で否定・五肢・組合せ表を確認。後年値で置換しない。\par

\paragraph{問われる核心} 2015年策定の2030電源見通しを読む。

\paragraph{必要知識} 策定年と目標年、電源構成と自給率、見通し。

\paragraph{考え方・正答を選ぶ理由（査読待ち初稿）} 2015の見通しの2030構成は、25–27\%がLNGで石油は約3\%。従って1が不適切。再エネ22–24\%、原子力20–22\%と、需給の省エネ・自給率目標は当時見通しの範囲。2030を対象とする後年改訂計画の比率をこの問題へ遡及しない。

\Needspace{4\baselineskip}\paragraph{選択肢の検討（査読待ち初稿）}\mbox{}\par

\noindent\textbf{1：不適切・選択対象} 石油25–27\%。判定と一致する。\par

\noindent\textbf{2：条件内で適切} 再エネ22–24\%。式・定義・対応の独立判定と一致しない。\par

\noindent\textbf{3：条件内で適切} 原子力20–22\%。式・定義・対応の独立判定と一致しない。\par

\noindent\textbf{4：条件内で適切} 自給率約25\%と準国産。式・定義・対応の独立判定と一致しない。\par

\noindent\textbf{5：条件内で適切} 省エネで最終消費低下。式・定義・対応の独立判定と一致しない。\par

\paragraph{再利用する考え方} 策定年と目標年、電源構成と自給率、見通し。 指定年度・対象範囲・単位・否定を先に固定し、他年度の番号や現行値を転用しない。

\paragraph{共通理論} \hyperref[theory:env:energymix2015]{第\ref*{theory:env:energymix2015}節}

\par 参考：\cite{env58_mix2015}。

\par 難易度 3（校準初稿）、優先度未評価。

\section{専門職責任は職務規定だけでは尽くせない}

\subsection{平成26年度 問5-5}
\label{q:h26-reg-basic-5-5}

\noindent\textbf{原典：} sources/H26/basic.pdf、PDF物理17頁。

\textbf{公式正答：} 3（single）。

\paragraph{問題の要旨} 専門職責任は職務規定だけでは尽くせない。

\Needspace{4\baselineskip}\paragraph{確認した条件・記述（独自要約）}\mbox{}\par

\noindent\textbf{時点：} 平成26年度に指定された資料・制度。\par

\noindent\textbf{照合：} 原典21画像で否定・五肢・組合せ表を確認。後年値で置換しない。\par

\paragraph{問われる核心} 専門職責任は職務規定だけでは尽くせない。

\paragraph{必要知識} 力量、公益、説明責任、法と倫理の範囲。

\paragraph{考え方・正答を選ぶ理由（査読待ち初稿）} 規定に明示がないことを責任免除の根拠とする3が不適切。専門職は力量の範囲、安全・公益、説明と結果への責任を考える。依頼業務が力量を超える場合、能力ある者との協働や適切な引継ぎを含め、無能力のまま責任を引き受けることを避ける。

\Needspace{4\baselineskip}\paragraph{選択肢の検討（査読待ち初稿）}\mbox{}\par

\noindent\textbf{1：条件内で適切} 職業で規範が異なる場合。式・定義・対応の独立判定と一致しない。\par

\noindent\textbf{2：条件内で適切} 力量外を引受けない。式・定義・対応の独立判定と一致しない。\par

\noindent\textbf{3：不適切・選択対象} 明文規定なしなら責任なし。判定と一致する。\par

\noindent\textbf{4：条件内で適切} 行動規範は変化。式・定義・対応の独立判定と一致しない。\par

\noindent\textbf{5：条件内で適切} 専門知識は高い基準。式・定義・対応の独立判定と一致しない。\par

\paragraph{再利用する考え方} 力量、公益、説明責任、法と倫理の範囲。 指定年度・対象範囲・単位・否定を先に固定し、他年度の番号や現行値を転用しない。

\paragraph{共通理論} \hyperref[theory:env:ethics]{第\ref*{theory:env:ethics}節}

\par 参考：\cite{ipej_ethics2011}。

\par 難易度 2（校準初稿）、優先度未評価。

\section{2012開始のFITの原則と拒否事由}

\subsection{平成26年度 問5-3}
\label{q:h26-reg-basic-5-3}

\noindent\textbf{原典：} sources/H26/basic.pdf、PDF物理17頁。

\textbf{公式正答：} 5（single）。

\par\textbf{検討未確定：} 2013METI資料の検索本文で正当拒否事由を確認したがPDF全文取得は失敗。出題時価格決定と対象水力条件の詳細、法令履歴は追加照合。

\paragraph{問題の要旨} 2012開始のFITの原則と拒否事由。

\Needspace{4\baselineskip}\paragraph{確認した条件・記述（独自要約）}\mbox{}\par

\noindent\textbf{時点：} 平成26年度に指定された資料・制度。\par

\noindent\textbf{照合：} 原典21画像で否定・五肢・組合せ表を確認。後年値で置換しない。\par

\paragraph{問われる核心} 2012開始のFITの原則と拒否事由。

\paragraph{必要知識} 固定価格、買取と接続、賦課金、例外。

\paragraph{考え方・正答を選ぶ理由（査読待ち初稿）} 正当な理由がある場合まで接続拒否を一切できないとする5は誤り。2012開始の対象は太陽光・風力・水力・地熱・バイオマスで、調達価格の設定、使用量に応じた賦課、費用調整という仕組みを区別する。例外を削って無条件義務にしない。

\Needspace{4\baselineskip}\paragraph{選択肢の検討（査読待ち初稿）}\mbox{}\par

\noindent\textbf{1：条件内で適切} 五種類の再エネ。式・定義・対応の独立判定と一致しない。\par

\noindent\textbf{2：条件内で適切} 大臣が年度価格設定。式・定義・対応の独立判定と一致しない。\par

\noindent\textbf{3：条件内で適切} 電力量に応じ賦課。式・定義・対応の独立判定と一致しない。\par

\noindent\textbf{4：条件内で適切} 地域負担差を調整。式・定義・対応の独立判定と一致しない。\par

\noindent\textbf{5：不適切・選択対象} 接続は一切拒否不可。判定と一致する。\par

\paragraph{再利用する考え方} 固定価格、買取と接続、賦課金、例外。 指定年度・対象範囲・単位・否定を先に固定し、他年度の番号や現行値を転用しない。

\paragraph{共通理論} \hyperref[theory:env:fit]{第\ref*{theory:env:fit}節}

\par 参考：\cite{env58_fit2013}。

\par 難易度 3（校準初稿）、優先度未評価。

\section{複合発電の高温側から低温側への順}

\subsection{平成28年度 問5-3}
\label{q:h28-reg-basic-5-3}

\noindent\textbf{原典：} sources/H28/basic.pdf、PDF物理22頁。

\textbf{公式正答：} 3（single）。

\paragraph{問題の要旨} 複合発電の高温側から低温側への順。

\Needspace{4\baselineskip}\paragraph{確認した条件・記述（独自要約）}\mbox{}\par

\noindent\textbf{時点：} 平成28年度に指定された資料・制度。\par

\noindent\textbf{照合：} 原典21画像で否定・五肢・組合せ表を確認。後年値で置換しない。\par

\paragraph{問われる核心} 複合発電の高温側から低温側への順。

\paragraph{必要知識} 温度段階、電気化学、静電蓄積、エネルギー変換。

\paragraph{考え方・正答を選ぶ理由（査読待ち初稿）} コンバインドは燃焼ガスでガスタービンを回し、その高温排熱から蒸気を作って蒸気タービンへ進む。逆順の3が不適切。二次電池は再充電、EDLCは界面で蓄える。揚水は発電を消費側から切り離して電気を水の位置エネルギーへ変換する貯蔵で、発電時だけの効率を往復効率としない。

\Needspace{4\baselineskip}\paragraph{選択肢の検討（査読待ち初稿）}\mbox{}\par

\noindent\textbf{1：条件内で適切} 一次使切・二次充電。式・定義・対応の独立判定と一致しない。\par

\noindent\textbf{2：条件内で適切} EDLC急速・長サイクル。式・定義・対応の独立判定と一致しない。\par

\noindent\textbf{3：不適切・選択対象} 先に蒸気タービン次に廃熱ガスタービン。判定と一致する。\par

\noindent\textbf{4：条件内で適切} 燃料電池は水電解の逆。式・定義・対応の独立判定と一致しない。\par

\noindent\textbf{5：条件内で適切} 揚水は位置エネルギーに貯蔵。式・定義・対応の独立判定と一致しない。\par

\paragraph{再利用する考え方} 温度段階、電気化学、静電蓄積、エネルギー変換。 指定年度・対象範囲・単位・否定を先に固定し、他年度の番号や現行値を転用しない。

\paragraph{共通理論} \hyperref[theory:env:generation]{第\ref*{theory:env:generation}節}

\par 難易度 2（校準初稿）、優先度未評価。

\section{発熱量は指定表と単位質量で比較}

\subsection{平成24年度 問5-2}
\label{q:h24-reg-basic-5-2}

\noindent\textbf{原典：} sources/H24/basic.pdf、PDF物理15頁。

\textbf{公式正答：} 4（single）。

\par\textbf{検討未確定：} 2005年度指定表の全四値を同じ定義で取得する作業は未完了。2023改訂資料の歴史表を補助参照し、後年値へ置換しない。

\paragraph{問題の要旨} 発熱量は指定表と単位質量で比較。

\Needspace{4\baselineskip}\paragraph{確認した条件・記述（独自要約）}\mbox{}\par

\noindent\textbf{時点：} 平成24年度に指定された資料・制度。\par

\noindent\textbf{照合：} 原典21画像で否定・五肢・組合せ表を確認。後年値で置換しない。\par

\paragraph{問われる核心} 発熱量は指定表と単位質量で比較。

\paragraph{必要知識} 高位・低位発熱量、単位、燃料品種、密度換算。

\paragraph{考え方・正答を選ぶ理由（査読待ち初稿）} 指定された燃料種の質量当たり発熱量は天然ガスCが最大、原油A、輸入原料炭B、廃材Dの順なので4。天然ガスの表が体積単位なら密度でMJ/kgへ換算する。液体のMJ/Lを質量当たり値へ直接並べない。

\Needspace{4\baselineskip}\paragraph{選択肢の検討（査読待ち初稿）}\mbox{}\par

\noindent\textbf{1：不適切} A>B>C>D。式・定義・対応の独立判定と一致しない。\par

\noindent\textbf{2：不適切} A>B>D>C。式・定義・対応の独立判定と一致しない。\par

\noindent\textbf{3：不適切} B>C>D>A。式・定義・対応の独立判定と一致しない。\par

\noindent\textbf{4：適切・選択対象} C>A>B>D。判定と一致する。\par

\noindent\textbf{5：不適切} C>A>D>B。式・定義・対応の独立判定と一致しない。\par

\paragraph{再利用する考え方} 高位・低位発熱量、単位、燃料品種、密度換算。 指定年度・対象範囲・単位・否定を先に固定し、他年度の番号や現行値を転用しない。

\paragraph{共通理論} \hyperref[theory:env:heating]{第\ref*{theory:env:heating}節}

\par 参考：\cite{env58_heating}。

\par 難易度 3（校準初稿）、優先度未評価。

\section{発表・発明を古い順へ並べる}

\subsection{平成25年度 問5-6}
\label{q:h25-reg-basic-5-6}

\noindent\textbf{原典：} sources/H25/basic.pdf、PDF物理22頁。

\textbf{公式正答：} 1（single）。

\par\textbf{検討未確定：} 全五業績の初出一次史料は未取得。特にハレー彗星の「発見」は周期性の同定を指す出題表現として留保。

\paragraph{問題の要旨} 発表・発明を古い順へ並べる。

\Needspace{4\baselineskip}\paragraph{確認した条件・記述（独自要約）}\mbox{}\par

\noindent\textbf{時点：} 平成25年度に指定された資料・制度。\par

\noindent\textbf{照合：} 原典21画像で否定・五肢・組合せ表を確認。後年値で置換しない。\par

\paragraph{問われる核心} 発表・発明を古い順へ並べる。

\paragraph{必要知識} 科学史の世紀区分、業績の段階、順序制約。

\paragraph{考え方・正答を選ぶ理由（査読待ち初稿）} ハレーの周期彗星の予測・同定が18世紀初頭、Wattの蒸気機関改良が18世紀後半、Darwin–Wallaceの自然選択発表が1858、一般相対論が1915、Fukuiのフロンティア理論が1952。エ→ア→イ→オ→ウで1。発見と理論発表、特許と実用化のどの段階をいうかも確認する。

\Needspace{4\baselineskip}\paragraph{選択肢の検討（査読待ち初稿）}\mbox{}\par

\noindent\textbf{1：適切・選択対象} エア イオウ。判定と一致する。\par

\noindent\textbf{2：不適切} アイエウオ。式・定義・対応の独立判定と一致しない。\par

\noindent\textbf{3：不適切} エアイウオ。式・定義・対応の独立判定と一致しない。\par

\noindent\textbf{4：不適切} アエオイウ。式・定義・対応の独立判定と一致しない。\par

\noindent\textbf{5：不適切} アイエオウ。式・定義・対応の独立判定と一致しない。\par

\paragraph{再利用する考え方} 科学史の世紀区分、業績の段階、順序制約。 指定年度・対象範囲・単位・否定を先に固定し、他年度の番号や現行値を転用しない。

\paragraph{共通理論} \hyperref[theory:env:historyorder]{第\ref*{theory:env:historyorder}節}

\par 参考：\cite{env58_steam}。

\par 難易度 3（校準初稿）、優先度未評価。

\subsection{平成26年度 問5-6}
\label{q:h26-reg-basic-5-6}

\noindent\textbf{原典：} sources/H26/basic.pdf、PDF物理18頁。

\textbf{公式正答：} 4（single）。

\par\textbf{検討未確定：} Galileo・Wilsonの初出史料は追加照合。ACSの1935生産選定と1939商業化を分ける。

\paragraph{問題の要旨} 観測・動力・検出・高分子の技術史。

\Needspace{4\baselineskip}\paragraph{確認した条件・記述（独自要約）}\mbox{}\par

\noindent\textbf{時点：} 平成26年度に指定された資料・制度。\par

\noindent\textbf{照合：} 原典21画像で否定・五肢・組合せ表を確認。後年値で置換しない。\par

\paragraph{問われる核心} 観測・動力・検出・高分子の技術史。

\paragraph{必要知識} 改良と新機構、検出技術、合成と商業化。

\paragraph{考え方・正答を選ぶ理由（査読待ち初稿）} Galileoの望遠鏡観測が1609年頃、Newcomen機関1712、Wattの改良特許1769、Wilson霧箱は20世紀初頭、Carothersのナイロンは1930年代。ウ→ア→イ→オ→エの4。ナイロンの合成・生産選定と1939の商業化を同一年にせず、霧箱より後という比較に必要な範囲で扱う。

\Needspace{4\baselineskip}\paragraph{選択肢の検討（査読待ち初稿）}\mbox{}\par

\noindent\textbf{1：不適切} イアウオエ。式・定義・対応の独立判定と一致しない。\par

\noindent\textbf{2：不適切} イウアオエ。式・定義・対応の独立判定と一致しない。\par

\noindent\textbf{3：不適切} イアウエオ。式・定義・対応の独立判定と一致しない。\par

\noindent\textbf{4：適切・選択対象} ウアイオエ。判定と一致する。\par

\noindent\textbf{5：不適切} ウイアエオ。式・定義・対応の独立判定と一致しない。\par

\paragraph{再利用する考え方} 改良と新機構、検出技術、合成と商業化。 指定年度・対象範囲・単位・否定を先に固定し、他年度の番号や現行値を転用しない。

\paragraph{共通理論} \hyperref[theory:env:historyorder]{第\ref*{theory:env:historyorder}節}

\par 参考：\cite{env58_steam,env58_nylon}。

\par 難易度 3（校準初稿）、優先度未評価。

\subsection{平成28年度 問5-6}
\label{q:h28-reg-basic-5-6}

\noindent\textbf{原典：} sources/H28/basic.pdf、PDF物理24頁。

\textbf{公式正答：} 4（single）。

\par\textbf{検討未確定：} 雷の実験・三極管・トランジスタ各原報と特許の一次史料は追加照合。三人の貢献を同一作業に単純化しない。

\paragraph{問題の要旨} 電気・化学・電子の代表的業績の順序。

\Needspace{4\baselineskip}\paragraph{確認した条件・記述（独自要約）}\mbox{}\par

\noindent\textbf{時点：} 平成28年度に指定された資料・制度。\par

\noindent\textbf{照合：} 原典21画像で否定・五肢・組合せ表を確認。後年値で置換しない。\par

\paragraph{問われる核心} 電気・化学・電子の代表的業績の順序。

\paragraph{必要知識} 周期律、放射性元素、電子管と半導体、年代順。

\paragraph{考え方・正答を選ぶ理由（査読待ち初稿）} 雷の電気的性質のFranklinが18世紀、Mendeleevの周期律1869、Curie夫妻の新元素1898、De Forestの三極管1906頃、トランジスタ1947。ウ→ア→イ→オ→エの4。電子管から半導体素子への世紀をまたぐ発展と、発見・発明・商用普及の段階を分ける。

\Needspace{4\baselineskip}\paragraph{選択肢の検討（査読待ち初稿）}\mbox{}\par

\noindent\textbf{1：不適切} アイウエオ。式・定義・対応の独立判定と一致しない。\par

\noindent\textbf{2：不適切} アイウオエ。式・定義・対応の独立判定と一致しない。\par

\noindent\textbf{3：不適切} ウイアオエ。式・定義・対応の独立判定と一致しない。\par

\noindent\textbf{4：適切・選択対象} ウアイオエ。判定と一致する。\par

\noindent\textbf{5：不適切} ウイアエオ。式・定義・対応の独立判定と一致しない。\par

\paragraph{再利用する考え方} 周期律、放射性元素、電子管と半導体、年代順。 指定年度・対象範囲・単位・否定を先に固定し、他年度の番号や現行値を転用しない。

\paragraph{共通理論} \hyperref[theory:env:historyorder]{第\ref*{theory:env:historyorder}節}

\par 参考：\cite{env58_curie}。

\par 難易度 3（校準初稿）、優先度未評価。

\section{人物と業績：放射能とラジウムを区別}

\subsection{平成23年度 問5-5}
\label{q:h23-reg-basic-5-5}

\noindent\textbf{原典：} sources/H23/basic.pdf、PDF物理17頁。

\textbf{公式正答：} 4（single）。

\par\textbf{検討未確定：} 原典は「アントワーヌ・ベクレル」と記載し人名が曖昧。放射能発見者Henri Becquerelとの対応を区別。Eastmanの「フィルム乾板」も乾板とロールフィルムをまとめた表現であり、原典を厳密な発明年表に一般化しない。

\paragraph{問題の要旨} 人物と業績：放射能とラジウムを区別。

\Needspace{4\baselineskip}\paragraph{確認した条件・記述（独自要約）}\mbox{}\par

\noindent\textbf{時点：} 平成23年度に指定された資料・制度。\par

\noindent\textbf{照合：} 原典21画像で否定・五肢・組合せ表を確認。後年値で置換しない。\par

\paragraph{問われる核心} 人物と業績：放射能とラジウムを区別。

\paragraph{必要知識} 放射能現象と元素発見、発明と改良・普及の区別。

\paragraph{考え方・正答を選ぶ理由（査読待ち初稿）} ラジウムとポロニウムの発見は1898年のキュリー夫妻の仕事。Henri Becquerelは1896年ウランの自然放射を見いだしたので4が誤り。Galileoは望遠鏡を改良して天体観測へ用いた。Huygensの振子時計、Arkwrightの水力紡績、Eastmanの写真材料・乾板の工業化という他の対応は、出題の代表的業績の範囲で扱う。

\Needspace{4\baselineskip}\paragraph{選択肢の検討（査読待ち初稿）}\mbox{}\par

\noindent\textbf{1：条件内で適切} Galileo/天体望遠鏡で観測。望遠鏡の発明者そのものとする記述ではなく観測への利用。\par

\noindent\textbf{2：条件内で適切} Huygens/振子時計。振子を時間標準にする時計。\par

\noindent\textbf{3：条件内で適切} Arkwright/水力紡績機。産業革命の水力紡績機。\par

\noindent\textbf{4：不適切・選択対象} Becquerel/ラジウム発見。放射能の発見とラジウム発見を取り違える。\par

\noindent\textbf{5：条件内で適切} Eastman/写真用フィルム乾板。乾板と柔軟な写真フィルムの開発・商業化という原典の代表表現。\par

\paragraph{再利用する考え方} 放射能現象と元素発見、発明と改良・普及の区別。 指定年度・対象範囲・単位・否定を先に固定し、他年度の番号や現行値を転用しない。

\paragraph{共通理論} \hyperref[theory:env:historypairs]{第\ref*{theory:env:historypairs}節}

\par 参考：\cite{env58_curie}。

\par 難易度 3（校準初稿）、優先度未評価。

\subsection{平成27年度 問5-6}
\label{q:h27-reg-basic-5-6}

\noindent\textbf{原典：} sources/H27/basic.pdf、PDF物理18頁。

\textbf{公式正答：} 4（single）。

\par\textbf{検討未確定：} 原典は「アントワーヌ・ベクレル」と記載し人名が曖昧。放射能発見者Henri Becquerelとの対応を区別。Eastmanの「フィルム乾板」も乾板とロールフィルムをまとめた表現であり、原典を厳密な発明年表に一般化しない。

\paragraph{問題の要旨} 人物と業績：放射能とラジウムを区別。

\Needspace{4\baselineskip}\paragraph{確認した条件・記述（独自要約）}\mbox{}\par

\noindent\textbf{時点：} 平成27年度に指定された資料・制度。\par

\noindent\textbf{照合：} 原典21画像で否定・五肢・組合せ表を確認。後年値で置換しない。\par

\paragraph{問われる核心} 人物と業績：放射能とラジウムを区別。

\paragraph{必要知識} 放射能現象と元素発見、発明と改良・普及の区別。

\paragraph{考え方・正答を選ぶ理由（査読待ち初稿）} ラジウムとポロニウムの発見は1898年のキュリー夫妻の仕事。Henri Becquerelは1896年ウランの自然放射を見いだしたので4が誤り。Galileoは望遠鏡を改良して天体観測へ用いた。Huygensの振子時計、Arkwrightの水力紡績、Eastmanの写真材料・乾板の工業化という他の対応は、出題の代表的業績の範囲で扱う。

\Needspace{4\baselineskip}\paragraph{選択肢の検討（査読待ち初稿）}\mbox{}\par

\noindent\textbf{1：条件内で適切} Galileo/天体望遠鏡で観測。望遠鏡の発明者そのものとする記述ではなく観測への利用。\par

\noindent\textbf{2：条件内で適切} Huygens/振子時計。振子を時間標準にする時計。\par

\noindent\textbf{3：条件内で適切} Arkwright/水力紡績機。産業革命の水力紡績機。\par

\noindent\textbf{4：不適切・選択対象} Becquerel/ラジウム発見。放射能の発見とラジウム発見を取り違える。\par

\noindent\textbf{5：条件内で適切} Eastman/写真用フィルム乾板。乾板と柔軟な写真フィルムの開発・商業化という原典の代表表現。\par

\paragraph{再利用する考え方} 放射能現象と元素発見、発明と改良・普及の区別。 指定年度・対象範囲・単位・否定を先に固定し、他年度の番号や現行値を転用しない。

\paragraph{共通理論} \hyperref[theory:env:historypairs]{第\ref*{theory:env:historypairs}節}

\par 参考：\cite{env58_curie}。

\par 難易度 3（校準初稿）、優先度未評価。

\section{産業革命の改良・工場化と担い手}

\subsection{平成24年度 問5-5}
\label{q:h24-reg-basic-5-5}

\noindent\textbf{原典：} sources/H24/basic.pdf、PDF物理17頁。

\textbf{公式正答：} 4（single）。

\par\textbf{検討未確定：} からくりからの転用という肢3の個別史料は追加照合。一般論を特定発明の確定年表に拡大しない。

\paragraph{問題の要旨} 産業革命の改良・工場化と担い手。

\Needspace{4\baselineskip}\paragraph{確認した条件・記述（独自要約）}\mbox{}\par

\noindent\textbf{時点：} 平成24年度に指定された資料・制度。\par

\noindent\textbf{照合：} 原典21画像で否定・五肢・組合せ表を確認。後年値で置換しない。\par

\paragraph{問われる核心} 産業革命の改良・工場化と担い手。

\paragraph{必要知識} 発明、改良、普及、研究組織と生産制度。

\paragraph{考え方・正答を選ぶ理由（査読待ち初稿）} 産業革命期の代表的機械は職人・発明家・事業者の改良と事業化で広まり、当時の欧州大学研究者が開発を主導した例が多いという4は不適切。新規性だけでなく既存機構の組合せ、工場制度、生産と社会への影響も見る。Wattと大学との接点のような個別事例は、全体を大学主導とする証明ではない。

\Needspace{4\baselineskip}\paragraph{選択肢の検討（査読待ち初稿）}\mbox{}\par

\noindent\textbf{1：条件内で適切} 既存技術の改良の積重ね。式・定義・対応の独立判定と一致しない。\par

\noindent\textbf{2：条件内で適切} 機械化で大工場へ。式・定義・対応の独立判定と一致しない。\par

\noindent\textbf{3：条件内で適切} 娯楽自動機構から転用もある。式・定義・対応の独立判定と一致しない。\par

\noindent\textbf{4：不適切・選択対象} 大学研究者主導の産学協同例が多い。判定と一致する。\par

\noindent\textbf{5：条件内で適切} 労働と児童労働が社会問題。式・定義・対応の独立判定と一致しない。\par

\paragraph{再利用する考え方} 発明、改良、普及、研究組織と生産制度。 指定年度・対象範囲・単位・否定を先に固定し、他年度の番号や現行値を転用しない。

\paragraph{共通理論} \hyperref[theory:env:industrial]{第\ref*{theory:env:industrial}節}

\par 参考：\cite{env58_steam}。

\par 難易度 3（校準初稿）、優先度未評価。

\section{知的財産の範囲は特許発明より広い}

\subsection{平成27年度 問5-5}
\label{q:h27-reg-basic-5-5}

\noindent\textbf{原典：} sources/H27/basic.pdf、PDF物理17頁。

\textbf{公式正答：} 3（single）。

\paragraph{問題の要旨} 知的財産の範囲は特許発明より広い。

\Needspace{4\baselineskip}\paragraph{確認した条件・記述（独自要約）}\mbox{}\par

\noindent\textbf{時点：} 平成27年度に指定された資料・制度。\par

\noindent\textbf{照合：} 原典21画像で否定・五肢・組合せ表を確認。後年値で置換しない。\par

\paragraph{問われる核心} 知的財産の範囲は特許発明より広い。

\paragraph{必要知識} 発明、知的財産、産業財産権、基本法の責務。

\paragraph{考え方・正答を選ぶ理由（査読待ち初稿）} 特許法の発明は自然法則を利用した高度な技術的思想。知的財産基本法の知財は著作物・商標・商号・営業秘密なども含むので、自然法則利用へ限定する3が不適切。法律名とその定義の範囲を別々に確認する。

\Needspace{4\baselineskip}\paragraph{選択肢の検討（査読待ち初稿）}\mbox{}\par

\noindent\textbf{1：条件内で適切} 特許は発明保護利用と産業発達。式・定義・対応の独立判定と一致しない。\par

\noindent\textbf{2：条件内で適切} 発明は自然法則利用の高度技術思想。式・定義・対応の独立判定と一致しない。\par

\noindent\textbf{3：不適切・選択対象} 知財は自然法則利用のものだけ。判定と一致する。\par

\noindent\textbf{4：条件内で適切} 基本法は創造保護活用の基本事項。式・定義・対応の独立判定と一致しない。\par

\noindent\textbf{5：条件内で適切} 国に施策策定実施責務。式・定義・対応の独立判定と一致しない。\par

\paragraph{再利用する考え方} 発明、知的財産、産業財産権、基本法の責務。 指定年度・対象範囲・単位・否定を先に固定し、他年度の番号や現行値を転用しない。

\paragraph{共通理論} \hyperref[theory:env:ip]{第\ref*{theory:env:ip}節}

\par 参考：\cite{ip_patent_statute,ip_broad_scope}。

\par 難易度 2（校準初稿）、優先度未評価。

\section{気候緩和と適応、LCAの境界}

\subsection{平成25年度 問5-3}
\label{q:h25-reg-basic-5-3}

\noindent\textbf{原典：} sources/H25/basic.pdf、PDF物理20頁。

\textbf{公式正答：} 4（single）。

\par\textbf{検討未確定：} CDMとCFPの出題当時規則全文は未取得。定義の独立判断と条文・規格の確定は分ける。

\paragraph{問題の要旨} 気候緩和と適応、LCAの境界。

\Needspace{4\baselineskip}\paragraph{確認した条件・記述（独自要約）}\mbox{}\par

\noindent\textbf{時点：} 平成25年度に指定された資料・制度。\par

\noindent\textbf{照合：} 原典21画像で否定・五肢・組合せ表を確認。後年値で置換しない。\par

\paragraph{問われる核心} 気候緩和と適応、LCAの境界。

\paragraph{必要知識} 緩和・適応、機能単位とシステム境界、環境基準。

\paragraph{考え方・正答を選ぶ理由（査読待ち初稿）} 排出削減や吸収を通じて原因を減らすのが緩和。影響に対して被害を減らす対策が適応なので4が取り違え。CFPは温室効果ガスのCO2換算に着目し、LCAはより広い環境影響の評価。環境基準は環境の望ましい状態で、個別工場の排出許容値と同じではない。

\Needspace{4\baselineskip}\paragraph{選択肢の検討（査読待ち初稿）}\mbox{}\par

\noindent\textbf{1：条件内で適切} 環境基準は維持が望ましい。式・定義・対応の独立判定と一致しない。\par

\noindent\textbf{2：条件内で適切} CDMは京都の認証削減。式・定義・対応の独立判定と一致しない。\par

\noindent\textbf{3：条件内で適切} CFPは全過程CO2換算。式・定義・対応の独立判定と一致しない。\par

\noindent\textbf{4：不適切・選択対象} 適応は排出削減で原因抑制。判定と一致する。\par

\noindent\textbf{5：条件内で適切} LCAは資源から廃棄まで評価。式・定義・対応の独立判定と一致しない。\par

\paragraph{再利用する考え方} 緩和・適応、機能単位とシステム境界、環境基準。 指定年度・対象範囲・単位・否定を先に固定し、他年度の番号や現行値を転用しない。

\paragraph{共通理論} \hyperref[theory:env:lca]{第\ref*{theory:env:lca}節}

\par 参考：\cite{env58_pmstandard}。

\par 難易度 3（校準初稿）、優先度未評価。

\section{気液の同質量比較からLNGの体積比}

\subsection{平成23年度 問5-3}
\label{q:h23-reg-basic-5-3}

\noindent\textbf{原典：} sources/H23/basic.pdf、PDF物理16頁。

\textbf{公式正答：} 2（single）。

\paragraph{問題の要旨} 気液の同質量比較からLNGの体積比。

\Needspace{4\baselineskip}\paragraph{確認した条件・記述（独自要約）}\mbox{}\par

\noindent\textbf{気体：} 0℃1気圧、純CH4の理想気体。\par

\noindent\textbf{液体：} 密度425kg/m\textsuperscript{3}一定とする原典モデル。\par

\paragraph{問われる核心} 気液の同質量比較からLNGの体積比。

\paragraph{必要知識} 理想気体モル体積、分子量、密度と体積、単位。

\paragraph{考え方・正答を選ぶ理由（査読待ち初稿）} 同じ物質量を比べる。0℃1気圧で1molの理想気体は約22.4L、メタン質量16g。液体体積は0.016kg/(425kg/m\textsuperscript{3})=\allowbreak{}3.76×10\textasciicircum{}-5m\textsuperscript{3}、気体は.0224m\textsuperscript{3}。比は.00168≈1/595なので最も近い1/600。密度比ρgas/ρliquidでも同じ。冷却の温度条件を気体標準状態と混同しない。

\Needspace{4\baselineskip}\paragraph{選択肢の検討（査読待ち初稿）}\mbox{}\par

\noindent\textbf{1：不適切} 1/400。式・定義・対応の独立判定と一致しない。\par

\noindent\textbf{2：適切・選択対象} 1/600。判定と一致する。\par

\noindent\textbf{3：不適切} 1/800。式・定義・対応の独立判定と一致しない。\par

\noindent\textbf{4：不適切} 1/1000。式・定義・対応の独立判定と一致しない。\par

\noindent\textbf{5：不適切} 1/1200。式・定義・対応の独立判定と一致しない。\par

\paragraph{再利用する考え方} 理想気体モル体積、分子量、密度と体積、単位。 指定年度・対象範囲・単位・否定を先に固定し、他年度の番号や現行値を転用しない。

\paragraph{共通理論} \hyperref[theory:env:lng]{第\ref*{theory:env:lng}節}

\par 難易度 2（校準初稿）、優先度未評価。

\section{2012年石油統計の対象範囲を固定}

\subsection{平成25年度 問5-1}
\label{q:h25-reg-basic-5-1}

\noindent\textbf{原典：} sources/H25/basic.pdf、PDF物理19頁。

\textbf{公式正答：} 5（single）。

\par\textbf{検討未確定：} 2012暦年の財務省四項目原表を全て取得していない。取得白書は2013年度83.6\%とサウジ首位を支持する補助で、指定年の確定証跡とは区別する。

\paragraph{問題の要旨} 2012年石油統計の対象範囲を固定。

\Needspace{4\baselineskip}\paragraph{確認した条件・記述（独自要約）}\mbox{}\par

\noindent\textbf{時点：} 平成25年度に指定された資料・制度。\par

\noindent\textbf{照合：} 原典21画像で否定・五肢・組合せ表を確認。後年値で置換しない。\par

\paragraph{問われる核心} 2012年石油統計の対象範囲を固定。

\paragraph{必要知識} 暦年と年度、数量と金額、統計品目、輸送経路。

\paragraph{考え方・正答を選ぶ理由（査読待ち初稿）} 2012年の指定統計では中東依存度約83\%、主要経路ホルムズ海峡、最大輸入先サウジアラビア、原油と石油製品の総輸入額比約21\%で5。数量の依存度と金額比は別であり、石炭・LNGを含む鉱物性燃料全体の比と原油・製品だけの比も違う。

\Needspace{4\baselineskip}\paragraph{選択肢の検討（査読待ち初稿）}\mbox{}\par

\noindent\textbf{1：不適切} 93/マラッカ/クウェート/21。式・定義・対応の独立判定と一致しない。\par

\noindent\textbf{2：不適切} 93/ホルムズ/サウジ/42。式・定義・対応の独立判定と一致しない。\par

\noindent\textbf{3：不適切} 83/マラッカ/サウジ/42。式・定義・対応の独立判定と一致しない。\par

\noindent\textbf{4：不適切} 83/ホルムズ/クウェート/42。式・定義・対応の独立判定と一致しない。\par

\noindent\textbf{5：適切・選択対象} 83/ホルムズ/サウジ/21。判定と一致する。\par

\paragraph{再利用する考え方} 暦年と年度、数量と金額、統計品目、輸送経路。 指定年度・対象範囲・単位・否定を先に固定し、他年度の番号や現行値を転用しない。

\paragraph{共通理論} \hyperref[theory:env:oil2012]{第\ref*{theory:env:oil2012}節}

\par 参考：\cite{env58_oil}。

\par 難易度 3（校準初稿）、優先度未評価。

\section{リサイクル制度の対象と責任主体}

\subsection{平成24年度 問5-1}
\label{q:h24-reg-basic-5-1}

\noindent\textbf{原典：} sources/H24/basic.pdf、PDF物理15頁。

\textbf{公式正答：} 3（single）。

\par\textbf{検討未確定：} 出題時点の建設規模・バーゼル条約対象と例外は追加の法令履歴照合。現行FAQだけで全過去条文の確定とはしない。

\paragraph{問題の要旨} リサイクル制度の対象と責任主体。

\Needspace{4\baselineskip}\paragraph{確認した条件・記述（独自要約）}\mbox{}\par

\noindent\textbf{時点：} 平成24年度に指定された資料・制度。\par

\noindent\textbf{照合：} 原典21画像で否定・五肢・組合せ表を確認。後年値で置換しない。\par

\paragraph{問われる核心} リサイクル制度の対象と責任主体。

\paragraph{必要知識} 法の目的、対象条件、排出者・自治体・小売・製造者の役割。

\paragraph{考え方・正答を選ぶ理由（査読待ち初稿）} 3が家電リサイクル法の小売業者の役割に対応。基本法はまず発生抑制。容器包装は三品目のみではなく、分別収集と再商品化の役割・費用も異なる。建設は規模等の対象条件があり全工事・発注者への実施義務という全称が違う。バーゼルは再利用目的の有害廃棄物でも規制を受け、条約成立の背景の輸送方向も逆。

\Needspace{4\baselineskip}\paragraph{選択肢の検討（査読待ち初稿）}\mbox{}\par

\noindent\textbf{1：不適切} 基本法は再生利用最優先。式・定義・対応の独立判定と一致しない。\par

\noindent\textbf{2：不適切} 容器三品目の全費用を販売者。式・定義・対応の独立判定と一致しない。\par

\noindent\textbf{3：適切・選択対象} 家電小売の引取・製造者への引渡し。判定と一致する。\par

\noindent\textbf{4：不適切} 全建設工事の発注者が分別再資源化。式・定義・対応の独立判定と一致しない。\par

\noindent\textbf{5：不適切} 再利用なら有害物輸出規制なし。式・定義・対応の独立判定と一致しない。\par

\paragraph{再利用する考え方} 法の目的、対象条件、排出者・自治体・小売・製造者の役割。 指定年度・対象範囲・単位・否定を先に固定し、他年度の番号や現行値を転用しない。

\paragraph{共通理論} \hyperref[theory:env:recyclelaw]{第\ref*{theory:env:recyclelaw}節}

\par 参考：\cite{env_circular_law,env_appliance}。

\par 難易度 3（校準初稿）、優先度未評価。

\section{資源量と太陽光面積の桁を比較}

\subsection{平成26年度 問5-4}
\label{q:h26-reg-basic-5-4}

\noindent\textbf{原典：} sources/H26/basic.pdf、PDF物理17頁。

\textbf{公式正答：} 4（single）。

\par\textbf{検討未確定：} 概算の日射・需要は桁比較であり厳密な上限証明ではない。石炭石油の指定時点埋蔵量、海水U総量、全球純生産と人類エネルギーの各原表は追加照合。初回記録の昼夜連続1kW仮定を解説では日照条件を示す概算へ改めた。

\paragraph{問題の要旨} 資源量と太陽光面積の桁を比較。

\Needspace{4\baselineskip}\paragraph{確認した条件・記述（独自要約）}\mbox{}\par

\noindent\textbf{時点：} 平成26年度に指定された資料・制度。\par

\noindent\textbf{照合：} 原典21画像で否定・五肢・組合せ表を確認。後年値で置換しない。\par

\paragraph{問われる核心} 資源量と太陽光面積の桁を比較。

\paragraph{必要知識} 面積換算、電力とエネルギー、存在量と可採埋蔵量。

\paragraph{考え方・正答を選ぶ理由（査読待ち初稿）} 100km\textsuperscript{2}=10\textasciicircum{}8m\textsuperscript{2}。例えば1kW/m\textsuperscript{2}を昼間平均12h/日、効率を非現実的に100\%としても年約438TWhで、当時の日本の年間電力需要約1000TWhに届かない。現実には夜間・天候・入射角・変換損失でさらに小さいので4。資源の存在量と経済的に採れる確認埋蔵量も同じではなく、後者は探鉱・価格・技術で増え得る。

\Needspace{4\baselineskip}\paragraph{選択肢の検討（査読待ち初稿）}\mbox{}\par

\noindent\textbf{1：条件内で適切} 石炭埋蔵エネルギーが石油超。式・定義・対応の独立判定と一致しない。\par

\noindent\textbf{2：条件内で適切} 石油確認埋蔵量は危機時より増。式・定義・対応の独立判定と一致しない。\par

\noindent\textbf{3：条件内で適切} 海水Uは1億t超。式・定義・対応の独立判定と一致しない。\par

\noindent\textbf{4：不適切・選択対象} 100km\textsuperscript{2}太陽光が日本需要超。判定と一致する。\par

\noindent\textbf{5：条件内で適切} 全球の光合成固定が人類需要超。式・定義・対応の独立判定と一致しない。\par

\paragraph{再利用する考え方} 面積換算、電力とエネルギー、存在量と可採埋蔵量。 指定年度・対象範囲・単位・否定を先に固定し、他年度の番号や現行値を転用しない。

\paragraph{共通理論} \hyperref[theory:env:resource]{第\ref*{theory:env:resource}節}

\par 難易度 3（校準初稿）、優先度未評価。

\section{下水処理の段階と環境対策}

\subsection{平成25年度 問5-4}
\label{q:h25-reg-basic-5-4}

\noindent\textbf{原典：} sources/H25/basic.pdf、PDF物理21頁。

\textbf{公式正答：} 2（single）。

\par\textbf{検討未確定：} 当時土壌対策の掘削除去割合と施設の詳細法定仕様は追加照合。緑化の効果は気象・水・建物条件に依存。

\paragraph{問題の要旨} 下水処理の段階と環境対策。

\Needspace{4\baselineskip}\paragraph{確認した条件・記述（独自要約）}\mbox{}\par

\noindent\textbf{時点：} 平成25年度に指定された資料・制度。\par

\noindent\textbf{照合：} 原典21画像で否定・五肢・組合せ表を確認。後年値で置換しない。\par

\paragraph{問われる核心} 下水処理の段階と環境対策。

\paragraph{必要知識} 一次・二次・高度処理、物理と生物反応、汚染経路。

\paragraph{考え方・正答を選ぶ理由（査読待ち初稿）} 活性汚泥による生物処理は一般に二次処理。一次は沈殿等の物理処理であり、2が不適切。管理型処分場は浸出水等による汚染を防ぐ。緑化は遮蔽・蒸発散、ダイオキシン対策は燃焼と再合成防止の工程を分けて説明する。

\Needspace{4\baselineskip}\paragraph{選択肢の検討（査読待ち初稿）}\mbox{}\par

\noindent\textbf{1：条件内で適切} 管理型の遮水・浸出水処理。式・定義・対応の独立判定と一致しない。\par

\noindent\textbf{2：不適切・選択対象} 活性汚泥は一次処理。判定と一致する。\par

\noindent\textbf{3：条件内で適切} 緑化の表面温度と空調低減。式・定義・対応の独立判定と一致しない。\par

\noindent\textbf{4：条件内で適切} 土壌対策に掘削除去が多い。式・定義・対応の独立判定と一致しない。\par

\noindent\textbf{5：条件内で適切} 完全燃焼と排ガス急冷。式・定義・対応の独立判定と一致しない。\par

\paragraph{再利用する考え方} 一次・二次・高度処理、物理と生物反応、汚染経路。 指定年度・対象範囲・単位・否定を先に固定し、他年度の番号や現行値を転用しない。

\paragraph{共通理論} \hyperref[theory:env:treatment]{第\ref*{theory:env:treatment}節}

\par 難易度 3（校準初稿）、優先度未評価。

\section{温対法の排出抑制と吸収作用}

\subsection{平成27年度 問5-2}
\label{q:h27-reg-basic-5-2}

\noindent\textbf{原典：} sources/H27/basic.pdf、PDF物理15頁。

\textbf{公式正答：} 1（single）。

\par\textbf{検討未確定：} 2015施行時の全五条項を履歴条文で一括確認する作業は未完了。現行の2050目標を当時条文へ追加しない。

\paragraph{問題の要旨} 温対法の排出抑制と吸収作用。

\Needspace{4\baselineskip}\paragraph{確認した条件・記述（独自要約）}\mbox{}\par

\noindent\textbf{時点：} 平成27年度に指定された資料・制度。\par

\noindent\textbf{照合：} 原典21画像で否定・五肢・組合せ表を確認。後年値で置換しない。\par

\paragraph{問われる核心} 温対法の排出抑制と吸収作用。

\paragraph{必要知識} 目的規定、吸収と排出、政令対象、計画。

\paragraph{考え方・正答を選ぶ理由（査読待ち初稿）} 目的は排出抑制だけでなく吸収作用の保全・強化も含む。言及しないという1が誤り。対象ガス、算定報告の主体・規模、地方推進体制、計画見直しは別の制度事項なので、吸収と排出の二側面を軸に確認する。

\Needspace{4\baselineskip}\paragraph{選択肢の検討（査読待ち初稿）}\mbox{}\par

\noindent\textbf{1：不適切・選択対象} 森林吸収の保全に言及なし。判定と一致する。\par

\noindent\textbf{2：条件内で適切} 対象七種の温室効果ガス。式・定義・対応の独立判定と一致しない。\par

\noindent\textbf{3：条件内で適切} 一定規模事業所別の算定報告。式・定義・対応の独立判定と一致しない。\par

\noindent\textbf{4：条件内で適切} 知事が推進員委嘱とセンター指定。式・定義・対応の独立判定と一致しない。\par

\noindent\textbf{5：条件内で適切} 少なくとも三年ごと計画検討。式・定義・対応の独立判定と一致しない。\par

\paragraph{再利用する考え方} 目的規定、吸収と排出、政令対象、計画。 指定年度・対象範囲・単位・否定を先に固定し、他年度の番号や現行値を転用しない。

\paragraph{共通理論} \hyperref[theory:env:warminglaw]{第\ref*{theory:env:warminglaw}節}

\par 難易度 3（校準初稿）、優先度未評価。

\section{排出抑制から適正処分までの優先順}

\subsection{平成23年度 問5-2}
\label{q:h23-reg-basic-5-2}

\noindent\textbf{原典：} sources/H23/basic.pdf、PDF物理16頁。

\textbf{公式正答：} 2（single）。

\par\textbf{検討未確定：} 取得した2001資料は基本方針案と施策実績であり、2010白書の引用箇所全文の一致は追加照合。

\paragraph{問題の要旨} 排出抑制から適正処分までの優先順。

\Needspace{4\baselineskip}\paragraph{確認した条件・記述（独自要約）}\mbox{}\par

\noindent\textbf{時点：} 平成23年度に指定された資料・制度。\par

\noindent\textbf{照合：} 原典21画像で否定・五肢・組合せ表を確認。後年値で置換しない。\par

\paragraph{問われる核心} 排出抑制から適正処分までの優先順。

\paragraph{必要知識} 3R、循環利用と処分、基準年度と目標年度。

\paragraph{考え方・正答を選ぶ理由（査読待ち初稿）} 旧来の廃棄物行政の公衆衛生と公害対策に、循環型社会形成が加わる。優先は排出抑制→再使用→再生利用→熱回収→残るものの適正処分。2001基本方針の最終処分目標は1997基準の概ね半分なので2。処分を埋立の一方式だけへ狭めない。

\Needspace{4\baselineskip}\paragraph{選択肢の検討（査読待ち初稿）}\mbox{}\par

\noindent\textbf{1：不適切} 生活環境/再資源化/適正処理/半分。式・定義・対応の独立判定と一致しない。\par

\noindent\textbf{2：適切・選択対象} 公衆衛生/再使用/適正処分/半分。判定と一致する。\par

\noindent\textbf{3：不適切} 生活環境/再資源化/適正処分/1/3。式・定義・対応の独立判定と一致しない。\par

\noindent\textbf{4：不適切} 公衆衛生/再資源化/適正処理/1/3。式・定義・対応の独立判定と一致しない。\par

\noindent\textbf{5：不適切} 生活環境/再使用/適正埋立/1/3。式・定義・対応の独立判定と一致しない。\par

\paragraph{再利用する考え方} 3R、循環利用と処分、基準年度と目標年度。 指定年度・対象範囲・単位・否定を先に固定し、他年度の番号や現行値を転用しない。

\paragraph{共通理論} \hyperref[theory:env:wasteorder]{第\ref*{theory:env:wasteorder}節}

\par 参考：\cite{env_circular_law,env58_wastegoal}。

\par 難易度 3（校準初稿）、優先度未評価。

\section{廃棄物用語の略称と対象}

\subsection{平成27年度 問5-1}
\label{q:h27-reg-basic-5-1}

\noindent\textbf{原典：} sources/H27/basic.pdf、PDF物理15頁。

\textbf{公式正答：} 2（single）。

\par\textbf{検討未確定：} 当時のRDF添加法・産廃業種条件・硫酸ピッチのガス生成に関する技術資料を追加照合。全RDFが同一添加処方という一般化をしない。

\paragraph{問題の要旨} 廃棄物用語の略称と対象。

\Needspace{4\baselineskip}\paragraph{確認した条件・記述（独自要約）}\mbox{}\par

\noindent\textbf{時点：} 平成27年度に指定された資料・制度。\par

\noindent\textbf{照合：} 原典21画像で否定・五肢・組合せ表を確認。後年値で置換しない。\par

\paragraph{問われる核心} 廃棄物用語の略称と対象。

\paragraph{必要知識} 略称、法定分類、処理特性と有害性。

\paragraph{考え方・正答を選ぶ理由（査読待ち初稿）} E-wasteはelectrical/electronic waste、廃電気電子機器。環境に良い又はほぼ全量再資源化可能という定義ではなく2が誤り。RDFは可燃廃物の加工燃料、バイオマスは化石資源を除く。産業廃棄物は事業活動由来なら全てという分類でなく法定種類と業種条件を確認する。

\Needspace{4\baselineskip}\paragraph{選択肢の検討（査読待ち初稿）}\mbox{}\par

\noindent\textbf{1：条件内で適切} RDFはごみ固形化燃料。式・定義・対応の独立判定と一致しない。\par

\noindent\textbf{2：不適切・選択対象} E-wasteはほぼ全て再資源化できるエコ廃物。判定と一致する。\par

\noindent\textbf{3：条件内で適切} バイオマスは化石除く生物資源。式・定義・対応の独立判定と一致しない。\par

\noindent\textbf{4：条件内で適切} 産廃は事業由来の法定20種類。式・定義・対応の独立判定と一致しない。\par

\noindent\textbf{5：条件内で適切} 硫酸ピッチは強酸性油泥。式・定義・対応の独立判定と一致しない。\par

\paragraph{再利用する考え方} 略称、法定分類、処理特性と有害性。 指定年度・対象範囲・単位・否定を先に固定し、他年度の番号や現行値を転用しない。

\paragraph{共通理論} \hyperref[theory:env:wasteterm]{第\ref*{theory:env:wasteterm}節}

\par 難易度 3（校準初稿）、優先度未評価。

\section{科学技術史・政策の経緯}

\subsection{平成29年度 問5-6}
\label{q:h29-reg-basic-5-6}

\noindent\textbf{原典：} sources/H29/basic.pdf、PDF物理20頁。

\textbf{公式正答：} 5（single）。

\par\textit{原典本文の校正と独自解説は未作成。分類のみの登録項目。}

\subsection{平成30年度 問5-5}
\label{q:h30-reg-basic-5-5}

\noindent\textbf{原典：} sources/H30/basic.pdf、PDF物理20頁。

\textbf{公式正答：} 3（single）。

\par\textit{原典本文の校正と独自解説は未作成。分類のみの登録項目。}

\subsection{令和元年度再試験 問5-5}
\label{q:r01-ret-basic-5-5}

\noindent\textbf{原典：} sources/R01\_retest/basic.pdf、PDF物理19頁。

\textbf{公式正答：} 4（single）。

\par\textit{原典本文の校正と独自解説は未作成。分類のみの登録項目。}

\subsection{令和2年度 問5-6}
\label{q:r02-reg-basic-5-6}

\noindent\textbf{原典：} sources/R02/basic.pdf、PDF物理22頁。

\textbf{公式正答：} 1（single）。

\par\textit{原典本文の校正と独自解説は未作成。分類のみの登録項目。}

\subsection{令和3年度 問5-5}
\label{q:r03-reg-basic-5-5}

\noindent\textbf{原典：} sources/R03/basic.pdf、PDF物理25頁。

\textbf{公式正答：} 4（single）。

\par\textit{原典本文の校正と独自解説は未作成。分類のみの登録項目。}

\subsection{令和3年度 問5-6}
\label{q:r03-reg-basic-5-6}

\noindent\textbf{原典：} sources/R03/basic.pdf、PDF物理26頁。

\textbf{公式正答：} 4（single）。

\par\textit{原典本文の校正と独自解説は未作成。分類のみの登録項目。}

\subsection{令和4年度 問5-6}
\label{q:r04-reg-basic-5-6}

\noindent\textbf{原典：} sources/R04/basic.pdf、PDF物理20頁。

\textbf{公式正答：} 3（single）。

\par\textit{原典本文の校正と独自解説は未作成。分類のみの登録項目。}

\subsection{令和5年度 問5-5}
\label{q:r05-reg-basic-5-5}

\noindent\textbf{原典：} sources/R05/basic.pdf、PDF物理18頁。

\textbf{公式正答：} 1（single）。

\par\textit{原典本文の校正と独自解説は未作成。分類のみの登録項目。}

\subsection{令和6年度 問5-6}
\label{q:r06-reg-basic-5-6}

\noindent\textbf{原典：} sources/R06/basic.pdf、PDF物理20頁。

\textbf{公式正答：} 2（single）。

\par\textit{原典本文の校正と独自解説は未作成。分類のみの登録項目。}

\subsection{令和7年度 問5-5}
\label{q:r07-reg-basic-5-5}

\noindent\textbf{原典：} sources/R07/basic.pdf、PDF物理18頁。

\textbf{公式正答：} 2（single）。

\par\textit{原典本文の校正と独自解説は未作成。分類のみの登録項目。}

\subsection{令和7年度 問5-6}
\label{q:r07-reg-basic-5-6}

\noindent\textbf{原典：} sources/R07/basic.pdf、PDF物理19頁。

\textbf{公式正答：} 4（single）。

\par\textit{原典本文の校正と独自解説は未作成。分類のみの登録項目。}

\section{独立正規変数の和差と標本平均}

\subsection{平成23年度 問5-4}
\label{q:h23-reg-basic-5-4}

\noindent\textbf{原典：} sources/H23/basic.pdf、PDF物理17頁。

\textbf{公式正答：} 3（single）。

\paragraph{問題の要旨} 独立な製品質量の標本平均と総質量について、平均と標準偏差を求める。

\Needspace{4\baselineskip}\paragraph{確認した条件・記述（独自要約）}\mbox{}\par

\noindent\textbf{判断方向：} 最も適切な記述又は組合せを選ぶ\par

\noindent\textbf{原典照合：} 14原典画像で16問の条件、数値、単位、表の列と全5候補を確認。\par

\noindent\textbf{モデル：} 独立性、正規性、一定発生率などの前提を個別に保持する。\par

\paragraph{問われる核心} 独立な製品質量の標本平均と総質量について、平均と標準偏差を求める。

\paragraph{必要知識} 独立変数の線形結合、平均の線形性、分散への係数の二乗、正規分布の閉性。

\paragraph{考え方・正答を選ぶ理由（査読待ち初稿）} 各質量を独立正規変数とする。Aの標本平均は \(\bar A=(A_1+\cdots+A_4)/4\) で、期待値80g、分散 \(4(4^2)/4^2=4\,\mathrm{g^2}\) である。総質量は \(T=\sum_{i=1}^4A_i+\sum_{i=1}^4B_i\) なので期待値 \(4(80)+4(120)=800\,\mathrm g\)、分散 \(4(4^2)+4(3^2)=100\,\mathrm{g^2}\) となる。標準偏差はそれぞれ2gと10gで、3を選ぶ。平均を取る操作と、四つの値を合計する操作を混同しない。

\Needspace{4\baselineskip}\paragraph{選択肢の検討（査読待ち初稿）}\mbox{}\par

\noindent\textbf{1：不適切} 平均20・偏差4・合計400・偏差10。Aの標本平均の期待値は80gであり、標本数で平均を割らない。総質量の期待値も800gである。\par

\noindent\textbf{2：不適切} 平均80・偏差4・合計800・偏差7。平均の偏差は4gではなく2gである。合計の偏差も標準偏差をそのまま足した7gではなく10gである。\par

\noindent\textbf{3：適切} 平均80・偏差2・合計800・偏差10。Aの4個平均の偏差は4/2=\allowbreak{}2g、合計の分散は4×16+4×9=\allowbreak{}100g\textsuperscript{2}なので10gである。\par

\noindent\textbf{4：不適切} 平均20・偏差2・合計400・偏差7。平均と総質量の期待値がともに違い、合計の偏差も違う。\par

\noindent\textbf{5：不適切} 平均80・偏差2・合計800・偏差7。平均部分は合うが、合計の標準偏差は7gではない。独立成分の分散を足して最後に平方根を取る。\par

\paragraph{再利用する考え方} 変数の意味と単位を先に定め、平均・分散・確率・判定を順に分ける。品質の要求、工程の監視、ロットの受入れを区別し、成立条件のない一般化を避ける。

\paragraph{共通理論} \hyperref[theory:quality:variance]{第\ref*{theory:quality:variance}節}, \hyperref[theory:quality:normal]{第\ref*{theory:quality:normal}節}, \hyperref[theory:quality:sample_mean]{第\ref*{theory:quality:sample_mean}節}

\par 参考：\cite{stat_nist_normal,stat_psu_linear}。

\paragraph{関連問題} \hyperref[q:r06-reg-basic-1-6]{R06-REG-BASIC-1-6}

\par 難易度 3（校準初稿）、優先度未評価。

\section{工程の品質保証と標準化}

\subsection{平成24年度 問5-4}
\label{q:h24-reg-basic-5-4}

\noindent\textbf{原典：} sources/H24/basic.pdf、PDF物理17頁。

\textbf{公式正答：} 2（single）。

\paragraph{問題の要旨} 最終検査と上流の品質保証を分け、プロセスの標準化と改善の関係を読む。

\Needspace{4\baselineskip}\paragraph{確認した条件・記述（独自要約）}\mbox{}\par

\noindent\textbf{判断方向：} 最も適切な記述又は組合せを選ぶ\par

\noindent\textbf{原典照合：} 14原典画像で16問の条件、数値、単位、表の列と全5候補を確認。\par

\noindent\textbf{モデル：} 独立性、正規性、一定発生率などの前提を個別に保持する。\par

\paragraph{問われる核心} 最終検査と上流の品質保証を分け、プロセスの標準化と改善の関係を読む。

\paragraph{必要知識} 検査と予防、上流工程、標準化、PDCA、プロセスと成果。

\paragraph{考え方・正答を選ぶ理由（査読待ち初稿）} 完成品を検査して不適合品を隔離することと、不適合を生む工程を上流から改善することを分ける。設問の文脈ではアが検査、イが源流管理、ウが標準化、エがプロセスであり2となる。標準を設定して実施し、結果を確認して更新する循環が品質保証を支える。標準化は永久に方法を固定することではなく、改善結果を次の標準へ反映することを含む。

\Needspace{4\baselineskip}\paragraph{選択肢の検討（査読待ち初稿）}\mbox{}\par

\noindent\textbf{1：不適切} 検査・三現管理・基準化・プロセス。現場・現物・現実という三現の確認と、上流で不適合原因を抑える源流管理は同じ説明ではない。\par

\noindent\textbf{2：適切} 検査・源流管理・標準化・プロセス。要求適合を判定する検査、上流からの源流管理、方法を定める標準化、プロセスを管理する流れでつながる。\par

\noindent\textbf{3：不適切} 検査・三現管理・基準化・アウトカム。三現管理・基準化の対応に加え、ここで管理する対象は一連のプロセスである。\par

\noindent\textbf{4：不適切} 評価・源流管理・標準化・アウトカム。評価は広い概念だが本問の要求適合判定は検査である。末尾も成果だけを意味するアウトカムではなくプロセスである。\par

\noindent\textbf{5：不適切} 評価・三現管理・標準化・プロセス。要求適合判定と上流の改善に対応する語が検査・源流管理になっていない。\par

\paragraph{再利用する考え方} 変数の意味と単位を先に定め、平均・分散・確率・判定を順に分ける。品質の要求、工程の監視、ロットの受入れを区別し、成立条件のない一般化を避ける。

\paragraph{共通理論} \hyperref[theory:quality:assurance]{第\ref*{theory:quality:assurance}節}, \hyperref[theory:quality:pdca]{第\ref*{theory:quality:pdca}節}

\par 参考：\cite{quality_process_jsqc,quality_iso_process}。

\paragraph{関連問題} \hyperref[q:r03-reg-basic-1-3]{R03-REG-BASIC-1-3}

\par 難易度 2（校準初稿）、優先度未評価。

\section{未確定}

\subsection{平成29年度 問5-2}
\label{q:h29-reg-basic-5-2}

\noindent\textbf{原典：} sources/H29/basic.pdf、PDF物理18頁。

\textbf{公式正答：} 1（single）。

\par\textit{原典本文の校正と独自解説は未作成。分類のみの登録項目。}

\subsection{平成29年度 問5-3}
\label{q:h29-reg-basic-5-3}

\noindent\textbf{原典：} sources/H29/basic.pdf、PDF物理19頁。

\textbf{公式正答：} 4（single）。

\par\textit{原典本文の校正と独自解説は未作成。分類のみの登録項目。}

\subsection{平成29年度 問5-4}
\label{q:h29-reg-basic-5-4}

\noindent\textbf{原典：} sources/H29/basic.pdf、PDF物理19頁。

\textbf{公式正答：} 2（single）。

\par\textit{原典本文の校正と独自解説は未作成。分類のみの登録項目。}

\subsection{平成30年度 問5-4}
\label{q:h30-reg-basic-5-4}

\noindent\textbf{原典：} sources/H30/basic.pdf、PDF物理19頁。

\textbf{公式正答：} 4（single）。

\par\textit{原典本文の校正と独自解説は未作成。分類のみの登録項目。}

\subsection{平成30年度 問5-6}
\label{q:h30-reg-basic-5-6}

\noindent\textbf{原典：} sources/H30/basic.pdf、PDF物理20頁。

\textbf{公式正答：} 5（single）。

\par\textit{原典本文の校正と独自解説は未作成。分類のみの登録項目。}

\subsection{令和元年度 問5-3}
\label{q:r01-reg-basic-5-3}

\noindent\textbf{原典：} sources/R01/basic.pdf、PDF物理20頁。

\textbf{公式正答：} 3（single）。

\par\textit{原典本文の校正と独自解説は未作成。分類のみの登録項目。}

\subsection{令和元年度 問5-4}
\label{q:r01-reg-basic-5-4}

\noindent\textbf{原典：} sources/R01/basic.pdf、PDF物理21頁。

\textbf{公式正答：} 4（single）。

\par\textit{原典本文の校正と独自解説は未作成。分類のみの登録項目。}

\subsection{令和元年度 問5-5}
\label{q:r01-reg-basic-5-5}

\noindent\textbf{原典：} sources/R01/basic.pdf、PDF物理21頁。

\textbf{公式正答：} 2（single）。

\par\textit{原典本文の校正と独自解説は未作成。分類のみの登録項目。}

\subsection{令和元年度再試験 問5-3}
\label{q:r01-ret-basic-5-3}

\noindent\textbf{原典：} sources/R01\_retest/basic.pdf、PDF物理18頁。

\textbf{公式正答：} 3（single）。

\par\textit{原典本文の校正と独自解説は未作成。分類のみの登録項目。}

\subsection{令和2年度 問5-2}
\label{q:r02-reg-basic-5-2}

\noindent\textbf{原典：} sources/R02/basic.pdf、PDF物理20頁。

\textbf{公式正答：} 3（single）。

\par\textit{原典本文の校正と独自解説は未作成。分類のみの登録項目。}

\subsection{令和2年度 問5-3}
\label{q:r02-reg-basic-5-3}

\noindent\textbf{原典：} sources/R02/basic.pdf、PDF物理20頁。

\textbf{公式正答：} 4（single）。

\par\textit{原典本文の校正と独自解説は未作成。分類のみの登録項目。}

\subsection{令和2年度 問5-4}
\label{q:r02-reg-basic-5-4}

\noindent\textbf{原典：} sources/R02/basic.pdf、PDF物理21頁。

\textbf{公式正答：} 1（single）。

\par\textit{原典本文の校正と独自解説は未作成。分類のみの登録項目。}

\subsection{令和2年度 問5-5}
\label{q:r02-reg-basic-5-5}

\noindent\textbf{原典：} sources/R02/basic.pdf、PDF物理21頁。

\textbf{公式正答：} 4|3（multiple\_accepted）。

\par\textit{原典本文の校正と独自解説は未作成。分類のみの登録項目。}

\subsection{令和3年度 問5-1}
\label{q:r03-reg-basic-5-1}

\noindent\textbf{原典：} sources/R03/basic.pdf、PDF物理21頁。

\textbf{公式正答：} 2（single）。

\par\textit{原典本文の校正と独自解説は未作成。分類のみの登録項目。}

\subsection{令和3年度 問5-4}
\label{q:r03-reg-basic-5-4}

\noindent\textbf{原典：} sources/R03/basic.pdf、PDF物理24頁。

\textbf{公式正答：} 2（single）。

\par\textit{原典本文の校正と独自解説は未作成。分類のみの登録項目。}

\subsection{令和4年度 問5-1}
\label{q:r04-reg-basic-5-1}

\noindent\textbf{原典：} sources/R04/basic.pdf、PDF物理17頁。

\textbf{公式正答：} 2（single）。

\par\textit{原典本文の校正と独自解説は未作成。分類のみの登録項目。}

\subsection{令和4年度 問5-4}
\label{q:r04-reg-basic-5-4}

\noindent\textbf{原典：} sources/R04/basic.pdf、PDF物理19頁。

\textbf{公式正答：} 5（single）。

\par\textit{原典本文の校正と独自解説は未作成。分類のみの登録項目。}

\subsection{令和5年度 問5-1}
\label{q:r05-reg-basic-5-1}

\noindent\textbf{原典：} sources/R05/basic.pdf、PDF物理16頁。

\textbf{公式正答：} 3（single）。

\par\textit{原典本文の校正と独自解説は未作成。分類のみの登録項目。}

\subsection{令和5年度 問5-2}
\label{q:r05-reg-basic-5-2}

\noindent\textbf{原典：} sources/R05/basic.pdf、PDF物理16頁。

\textbf{公式正答：} 5（single）。

\par\textit{原典本文の校正と独自解説は未作成。分類のみの登録項目。}

\subsection{令和5年度 問5-3}
\label{q:r05-reg-basic-5-3}

\noindent\textbf{原典：} sources/R05/basic.pdf、PDF物理17頁。

\textbf{公式正答：} 1（single）。

\par\textit{原典本文の校正と独自解説は未作成。分類のみの登録項目。}

\subsection{令和5年度 問5-4}
\label{q:r05-reg-basic-5-4}

\noindent\textbf{原典：} sources/R05/basic.pdf、PDF物理17頁。

\textbf{公式正答：} 2（single）。

\par\textit{原典本文の校正と独自解説は未作成。分類のみの登録項目。}

\subsection{令和5年度 問5-6}
\label{q:r05-reg-basic-5-6}

\noindent\textbf{原典：} sources/R05/basic.pdf、PDF物理18頁。

\textbf{公式正答：} 5（single）。

\par\textit{原典本文の校正と独自解説は未作成。分類のみの登録項目。}

\subsection{令和7年度 問5-3}
\label{q:r07-reg-basic-5-3}

\noindent\textbf{原典：} sources/R07/basic.pdf、PDF物理17頁。

\textbf{公式正答：} 4（single）。

\par\textit{原典本文の校正と独自解説は未作成。分類のみの登録項目。}
